% Uniform positive curvature in an open ball. A local quadratic schematic at left;
% right: straight segment from x* to a hypothetical second stationary point y stays in ball.
% Same figure supports strictness and isolatedness via Taylor and integral forms.
\begin{tikzpicture}[x=1cm,y=1cm,font=\small,>=Latex]
\begin{scope}
\node at (2.4,3.7) {Every nearby direction rises};
\draw[->,black] (0,.15)--(4.8,.15) node[right] {step};
\draw[thick,blue!70!black] plot[domain=.4:4.4,samples=80] (\x,{.55*(\x-2.4)^2+.55});
\draw[densely dashed] (0,.55)--(4.8,.55);
\fill (2.4,.55) circle (2pt) node[below] {$\mathbf{x}^*$};
\draw[<->,thick,orange!80!black] (3.55,.55)--(3.55,1.277) node[midway,right] {$\Delta f>0$};
\node at (2.4,-.5) {strictness: positive change in value};
\end{scope}
\begin{scope}[xshift=7cm]
\node at (2.1,3.7) {Positive curvature along the whole segment};
\draw[thick,blue!70!black,densely dashed] (2.1,1.7) circle (1.65);
\fill (2.1,1.7) circle (2pt) node[above] {$\mathbf{x}^*$};
\fill (3.15,1.15) circle (2pt) node[right] {$\mathbf{y}$};
\draw[->,thick,orange!80!black] (2.1,1.7)--(3.15,1.15);
\node[above left] at (2.85,1.38) {$\mathbf{p}$};
\node at (2.1,-.5) {isolation: gradients cannot both be zero};
\end{scope}
\end{tikzpicture}
